% !TeX root = ../tesi.tex

\chapter{Corrispondenza di Edelman-Greene}
\label{ch:edelman_greene}

Consideriamo ancora il gruppo simmetrico $S_n$ come generato dalle trasposizioni di
elementi adiacenti $s_1, s_2, \dots, s_{n-1}$, dove $s_i = (i , i+1)$.

Ricordiamo che una parola $v = v_1 v_2 \cdots v_k$ nell'alfabeto $\{1, 2, \dots, n-1 \}$
\`e associata alla permutazione $\w = s_{v_1} s_{v_2} \cdots s_{v_k}$ in $\sym_n$.

Inoltre, la parola $v$ si dice \emph{ridotta} se la sua lunghezza $k$ \`e la minima
possibile tra tutte le parole associate a $\w$, ovvero il numero di inversioni
di $\w$, che chiamiamo \emph{lunghezza} di $\w$ e denotiamo con $\len(\w)$.

Per ogni permutazione $\w \in \sym_n$, chiamo $\rw(\w)$ l'insieme delle 
parole ridotte di $\w$.

\begin{definition}
    Un tableau semistandard si dice \emph{ridotto} se $\rowreading(T)$ \`e una
    parola ridotta.
\end{definition}

Per ogni permutazione $\w \in \sym_n$, chiamo $\redtab(\w)$ l'insieme dei 
tableaux $T$ (ridotti) tali che $\rowreading(T) \in \rw(\w)$.
% TODO usa questa notazione

% TODO aggiungi esempio
TODO: aggiungi un esempio

\begin{remark}
    \label{rmk:tab_ridotto_crescenza}
    % NOTE ponderare il emph
    Un tableau ridotto, oltre che avere le colonne \emph{strettamente} crescenti, 
    in quanto semistandard, ha anche le righe \emph{strettamente} crescenti.
    Infatti, altrimenti il suo row-reading avrebbe due lettere consecutive.
\end{remark}
    

In questo capitolo descriviamo la corrispondenza di Edelman e Greene (EG), descritta
per la prima volta in \cite[Section~6]{edelman-greene-87}.

% TODO aggiungi cose su EG e RSK


\section{Algoritmo di inserimento EG}

Proprio come nella corrispondenza RSK, la corrispondenza EG \`e ottenuta iterando 
un algoritmo di inserimento, che descriviamo in questa sezione.

% L'algoritmo di inserimento EG associa a ogni parola ridotta $v$ un tableau 
% ridotto $P(v)$ tale che $v$ e $\rowreading(P(v))$ sono parole ridotte per lo stesso
% elemento nel gruppo simmetrico.

Siano $T$ un tableau ridotto e $k$ una lettera tale che $\rowreading(T)k$ sia una
parola ridotta per un qualche $\w \in \sym_n$.

Voglio definire $T \leftarrow k$, ovvero l'inserimento di $k$ dentro $T$, in modo
che $\rowreading(T \leftarrow k)$ sia una parola ridotta per lo stesso $\w$.

Guardiamo prima il caso in cui $T$ ha una sola riga e siano 
$i_1 < i_2 < \cdots < i_s$ le sue lettere.

% \[
%     T = \ytableaushort{ {i_1}  {i_2}  \cdots  {i_s}  }
% \]

Si possono verificare tre situazioni:
\begin{enumerate}
    \item \label{item:eg_ins1}
        se vale $i_s < k$ (o se $T=\varnothing$, ovvero $T$ \`e il tableau vuoto),
        allora aggiungo una casella con $k$ alla riga di $T$:
        \[
            \begin{ytableau}
                i_1 & i_2 & \cdots & i_s & k
            \end{ytableau};
        \]
    \item \label{item:eg_ins2} 
        altrimenti, se $k$ \emph{non} compare nella riga: sia $u$ il più piccolo
        intero positivo tale che $i_u > k$, allora $k$ prende il posto di $i_u$ e
        $i_u$ viene ``spinto'' (bumped) nella riga successiva. Il risultato è il seguente:
        \[
            \begin{ytableaubig}
i_u \\
                i_1 & i_2 & \cdots & k & \scriptstyle i_{u+1} & \cdots & i_s 
            \end{ytableaubig};
        \]
    \item \label{item:eg_ins3}
        altrimenti, se $k$ compare nella riga, necessariamente deve comparire anche 
        $k+1$, altrimenti $\rowreading(T)k$ non sarebbe una parola ridotta.
        In questo caso la riga rimane invariata, ma viene ``spinto'' nella riga successiva
        $k+1$. Il risultato \`e il seguente:
        \[
            \begin{ytableaubig}
                \scriptstyle k+1 \\
                i_1 & i_2 & \cdots & k & \scriptstyle k+1 & \cdots & i_s 
            \end{ytableaubig}.
        \]

\end{enumerate}

\begin{remark} \label{rmk:coxeter_righe}
    Almeno nel caso in cui $T$ ha una sola riga \`e facile vedere che 
    $\rowreading(T \leftarrow k)$ si ottiene da $\rowreading(T)k$ tramite operazioni
    di Coxeter:
    \begin{itemize}
        \item nel caso~\ref{item:eg_ins2}, $k$ commuta con tutti gli $i_j$ a destra di $i_u$,
            arrivando alla parola intermedia
            \[
                i_1 i_2 \cdots i_u k i_{u+1} \cdots i_s.
            \]
            Poi $i_u$ commuta con tutti gli $i_j$
            alla sua sinistra, in quanto tutti minori di $k$ ottenendo
            \[
                i_u i_1 \cdots i_{u-1} k i_{u+1} \cdots i_s.
            \]
        \item nel caso~\ref{item:eg_ins3}, come prima $k$ commuta con tutti gli $i_j$
            a destra di $k+1$, arrivando alla parola
            \[
                i_1 i_2 \cdots k (k+1) k \cdots i_s.
            \]
            Ora applico la braid relation per ottenere
            \[
                i_1 i_2 \cdots  (k+1) k (k+1) \cdots i_s.
            \]
            Infine $k+1$ commuta con tutti gli $i_j$ alla sua sinistra in quanto minori
            di $k$, ottenendo
            \[
                (k+1) i_1 i_2 \cdots k (k+1) \cdots i_s.
            \]
    \end{itemize}
\end{remark}

Nel caso in cui $T$ ha pi\`u di una riga, si ripete l'algoritmo per inserire $k$
nella prima riga di $T$, ma ora l'eventuale lettera ``spinta'' viene inserita
nella riga successiva con lo stesso algoritmo e cos\`i via fino a che non si verifica
il caso~\ref{item:eg_ins1}, ovvero l'inserimento
non ``spinge'' nessuna lettera, oppure la lettera ``spinta'' viene aggiunta a 
una riga precedentemente vuota.

% TODO ponderare se qui vogliamo mettere che sto implicitamente usando l'osservazione
% per dire che le cose giuste rimangono ridotte...

% L'Osservazione~\ref{rmk:coxeter_righe} ci garantisce che in ogni step di questo processo
% sia verificata l'ipotesi che 

\begin{remark}
    L'algoritmo di inserimento EG \`e identico a quello per l'inserimento RSK, 
    tranne per il caso in cui la lettera da inserire compare già nella riga.
    In quel caso l'unica differenza \`e che la riga in questione rimane invariata.
\end{remark}

\begin{example}
    \label{ex:inserimento_eg}
    Consideriamo per esempio $k = 3$ e $T$ come sotto:
    \[
        \ytableausetup{aligntableaux=center}
        T = 
        \ytableaunumber{4, 3 4 5, 1 2 4 6 }
        \ytableausetup{aligntableaux=bottom}.
    \]
    Si verifica che $\rowreading(T)k = 434512463$ \`e una parola ridotta per 
    $\w = s_4 s_3 s_4 s_5 s_1 s_2 s_4 s_6 s_3 = 2561473 \in \sym_7$.

    Calcoliamo ora $T \leftarrow k$:
    \[ 
        \ytableaunumber{4, 3 4 5, 1 2 4 6 \none {\none[\leftarrow 3]}}
        \qquad
        \ytableaunumber{4, 3 4 5 \none {\none[\leftarrow 4]}, 1 2 {*(green) 3} 6 }
        \qquad
        \ytableaunumber{4 \none {\none[\leftarrow 5]}, 3 {*(green) 4} 5 , 1 2 {*(green) 3} 6 }
        \qquad
        \ytableaunumber{4 {*(green) 5}, 3 {*(green) 4} 5 , 1 2 {*(green) 3} 6 }.
    \]
    
    Per ogni riga ho colorato la casella in cui si \`e inserita la nuova lettera
    (o quella in cui era gi\`a scritta).
    Queste caselle prendono il nome di \emph{percorso di inserimento}.
\end{example}



\begin{proposition}
    Siano $T$ un tableau ridotto e $k$ una lettera tale che $v = \rowreading(T)k$ sia una 
    parola ridotta per un qualche $\w \in \sym_n$.
    Allora:
    \begin{enumerate}
        \item il tableau $T \leftarrow k$ \`e semistandard;
        \item il percorso di inserimento si muove debolmente verso sinistra;
        \item il tableau $T \leftarrow k$ \`e ridotto 
            e $v' = \rowreading(T\leftarrow k)$ \`e una parola ridotta per lo stesso $\w$.
    \end{enumerate}
\end{proposition}

\begin{proof}
    Vediamo prima che $T \leftarrow k$ rimane semistandard.

    La crescenza delle righe segue immediatamente dalla definizione; verifichiamo la crescenza
    delle colonne.

    Immaginiamo di stare inserendo $a$ in una riga di $T$. Sia $b$ il primo elemento di 
    quella riga che \`e maggiore o uguale ad $a$. La situazione \`e la seguente:
    \[
        \ytableaushort 
        { \none \none {\scriptstyle >b}, 
        {\scriptstyle <a} {\scriptstyle <a} b \none \none \none {\none[\leftarrow a]}}
        * {4,5} 
        \quad;
    \]
    che, dopo l'inserimento, diventa:
    \[
        \ytableausetup{aligntableaux=center}
        \begin{array}{ll}
            \ytableaushort 
            { \none \none {\scriptstyle >a} \none \none {\none[\leftarrow a+1]}, 
            {\scriptstyle <a} {\scriptstyle <a} {*(green) a}  }
            * {4,5}
            &\text{ se $a=b$; } 
            \\[2em]
            \ytableaushort 
            { \none \none {\scriptstyle >b} \none \none {\none[\leftarrow b]}, 
            {\scriptstyle <a} {\scriptstyle <a} {*(green) a}  }
            * {4,5}
            &\text{ se $a < b$.}
        \end{array}
        \ytableausetup{aligntableaux=bottom}
    \]

    In entrambi i casi l'inserimento nella riga successiva avverr\`a in una casella
    debolmente a sinistra della casella contenente $a$, ovvero l'inserimento
    perserver\`a la crescenza delle colonne e il percorso di inserimento
    si muover\`a debolmente a sinistra.

    % Vediamo ora che $v' = \rowreading(T\leftarrow k)$ si ottiene da $v = \rowreading(T)k$
    % con una serie di relazioni di Coxeter, e quindi rappresentano lo stesso $\w$.
    % TODO
    TODO: dire qualcosa su questo punto, anche se segue rapidamente dall'osservazione.
\end{proof}

Enunciamo e dimostriamo due lemmi che saranno utili in seguito.

\begin{lemma} \label{lem:percorsi_strict}
    Sia $T$ un tableau ridotto e siano $a, b$ due lettere con $a < b$ e tali che 
    $\rowreading(T)ab$ sia una parola ridotta.
    Allora inserendo $a$ dentro $T$ e poi $b$ dentro $T \leftarrow a$, ovvero
    calcolando $T \leftarrow a \leftarrow b$, il percorso di inserimento di $b$ \`e 
    \emph{strettamente} a destra di quello di $a$.

    Inoltre, la casella aggiunta alla forma del tableau inserendo $b$ \`e
    strettamente a destra e debolmente sotto quella aggiunta inserendo $a$.
\end{lemma}

\begin{proof}
    % TODO
    TODO
\end{proof}

\begin{lemma} \label{lem:percorsi_weak}
    Sia $T$ un tableau ridotto e siano $a, b$ due lettere questa volta con $a > b$
    e tali che $\rowreading(T)ab$ sia una parola ridotta.
    Allora inserendo $a$ dentro $T$ e poi $b$ dentro $T \leftarrow a$, ovvero
    calcolando $T \leftarrow a \leftarrow b$, il percorso di inserimento di $b$ \`e 
    \emph{debolmente} a sinistra di quello di $a$.

    Inoltre, la casella aggiunta alla forma del tableau inserendo $b$ \`e
    debolmente a sinistra e strettamente sopra quella aggiunta inserendo $a$.
\end{lemma}

\begin{proof}
    % TODO
\end{proof}

% \begin{lemma} \label{lem:strip}
%     Nelle stesse ipotesi del lemma precedente, la casella aggiunta inserendo $b$ \`e
%     strettamente a destra e debolmente sotto quella aggiunta inserendo $a$.
% \end{lemma}
% 
% \begin{proof}
% \end{proof}

\section{Corrispondenza EG per parole ridotte}

Possiamo adesso definire il $P$-tableau della corrispondenza EG, ottenuto iterando 
l'algoritmo definito sopra.

\begin{definition}[$P$-tableau EG]
    Sia $\w$ una permutazione in $\sym_n$ e sia $v=v_1 v_2 \cdots v_s$ una sua
    parola ridotta. 
    Definisco $P(v)$, ovvero il $P$-tableau EG associato a $v$, come il tableau
    ridotto ottenuto partendo dal tableau vuoto e inserendo in successione 
    le lettere di $v$ da sinistra a destra.
    In simboli:
    \[
        P(v) = \varnothing \leftarrow v_1 \leftarrow v_2 \leftarrow
        \cdots \leftarrow v_s.
    \]
\end{definition}

\begin{example} \label{ex:p_tab1}
    Sia $v = 232123$, parola ridotta di $\w = s_2 s_3 s_2 s_1 s_2 s_3 = 4321 \in \sym_4$.
    Calcoliamo $P(v)$ inserendo una lettera alla volta:
    \[
        \ytableaunumber{2}
        \xrightarrow{\leftarrow 3}
        \ytableaunumber{2 3}
        \xrightarrow{\leftarrow 2}
        \ytableaunumber{3, 2 3}
        \xrightarrow{\leftarrow 1}
        \ytableaunumber{3, 2, 1 3}
        \xrightarrow{\leftarrow 2}
        \ytableaunumber{3, 2 3, 1 2}
        \xrightarrow{\leftarrow 3}
        \ytableaunumber{3, 2 3, 1 2 3}
        = P(v).
    \]

    Si verifica anche che $\rowreading(P(v)) = 323123$ \`e una parola ridotta per $\w$.

\end{example}

Dimostriamo alcune propriet\`a del $P$-tableau EG che useremo.

\begin{proposition}[Propriet\`a del $P$-tableau EG]
    La mappa $P$ di EG soddisfa le seguenti propriet\`a:
    \begin{enumerate}
        \item $P(v)$ \`e ridotto;
        \item $v$ e $\rowreading(P(v))$ sono parole ridotte per lo stesso elemento in $\sym_n$.
    \end{enumerate}
    % TODO
\end{proposition}

In completa analogia con la corrispondenza RSK, definiamo il $Q$-tableau EG, o
\emph{recording} tableau EG, che tiene traccia di dove vengono aggiunte le caselle 
e permette di invertire il processo.

\begin{definition}[$Q$-tableau EG]
    Sia $v = v_1 v_2 \cdots v_s$ una parola ridotta. 
    Definisco $Q(v)$, ovvero il $Q$-tableau EG associato a $v$, come il tableau di Young
    con la stessa forma di $P(v)$ nelle cui caselle \`e segnato l'ordine in cui sono state
    aggiunte durante la costruzione del $P$-tableau.

    Pi\`u formalmente, definisco la successione dei $P$-tableaux intermedi
    $P^{(0)}$, $P^{(1)}$, \dots, $P^{(s)}$, dove
    \begin{align*}
        P^{(0)} &= \varnothing, \\
        P^{(k)} &= P^{(k-1)} \leftarrow v_k \qquad \text{se $k>0$};
    \end{align*}
    da cui $P(v) = P^{(s)}$.
    Se $(i,j)$ \`e la casella di $P^{(k)} \setminus P^{(k-1)}$, allora scrivo $k$ nella 
    casella $(i,j)$ di $Q(v)$.
    
\end{definition}

\begin{example} \label{ex:q_tab1}
    Continuando l'Esempio~\ref{ex:p_tab1}, in cui avevamo scelto $v = 232123$,
    il $Q$-tableau \`e ottenuto nel seguente modo:
    \[
        \ytableaunumber{1}
        \longrightarrow
        \ytableaunumber{1 2}
        \longrightarrow
        \ytableaunumber{3, 1 2}
        \longrightarrow
        \ytableaunumber{4, 3, 1 2}
        \longrightarrow
        \ytableaunumber{4, 3 5, 1 2}
        \longrightarrow
        \ytableaunumber{4, 3 5, 1 2 6}
        =Q(v).
    \]
\end{example}

\begin{remark}
    Il $Q$-tableau \`e un tableau \emph{standard}.
\end{remark}

\begin{theorem}[{\cite[Theorem~6.25]{edelman-greene-87}}]
    \label{thm:eg1}
    Fissata una permutazione $\w \in \sym_n$, la corrispondenza EG definita da
    \[
        \varphi_{\mathrm{EG}} \colon v \mapsto (P(v), Q(v))
    \]
    \`e una bigezione tra:
    \begin{itemize}
        \item l'insieme $\rw(\w)$ delle parole ridotte di $\w$;
        \item l'insieme delle coppie $(P,Q)$ di tableaux di Young con la stessa forma,
            tali che $P$ sta in $\redtab(\w)$
            e $Q$ \`e standard.
    \end{itemize}
    
\end{theorem}

\begin{proof}
    Basta mostrare che, conoscendo la posizione dell'ultima casella
    inserita, l'algoritmo di inserimento EG pu\`o sempre essere invertito in modo univoco.
    % TODO dim EG

    Sia $k$ la lettera che \`e stata ``spinta'' che vogliamo riposizionare nella riga
    sottostante. 
    Si possono verificare due casi:
    \begin{itemize}
        \item se $k$ \emph{non} compare nella riga sottostante, allora $k$ era stata 
            ``spinta'' da un inserimento di tipo~\ref{item:eg_ins2} e quindi occupava la 
            casella della pi\`u grande lettera minore di $k$;
        \item se $k$ compare nella riga sottostante, allora deve comparire anche $k-1$
            (altrimenti il row-reading non sarebbe una parola ridotta), e quindi
            $k$ era stata ``spinta'' da un inserimento di tipo~\ref{item:eg_ins3}.
        \qedhere
    \end{itemize}
\end{proof}

\begin{example}
    % TODO forse mettere un esempio
\end{example}
    
Questa bigezione permette gi\`a di dare una risposta parziale alla domanda originale,
ovvero quante siano le parole ridotte di una permutazione $\w$.

\begin{definition}
    Per ogni permutazione $\w \in \sym_n$ e partizione 
    $\lambda \vdash \len(\w)$,
    definiamo $a_\w^\lambda$ come il numero di tableaux semistandard il cui row-reading
    \`e una parola ridotta per $\w$.
    In simboli:
    \[
        % \label{eq:def_a_w_lambda}
        a_\w^\lambda \coloneq \lvert \left\{ 
            T \in \redtab(\w) \mid \shape(T) = \lambda
        \right\} \rvert .
    \]

    Definiamo inoltre $f^\lambda$ come il numero di tableaux standard di forma $\lambda$.
\end{definition}

\begin{corollary}
    Per ogni permutazione $\w \in \sym_n$, vale
    \begin{equation}
        \lvert \rw(\w) \rvert = \sum_{\lambda} a_\w^\lambda f^\lambda,
    \end{equation}
    con $a_\w^\lambda$ e $f^\lambda$ come definiti subito sopra.
\end{corollary}

In generale, non \`e facile calcolare i valori di $a_\w^\lambda$.
In \cite{stanley-84}[Section~4], Stanley ha dato alcune condizioni su $\lambda$ 
necessarie affich\'e $a_\w^\lambda$ sia diverso da zero.

In alcuni casi particolari, invece, \`e possibile dare una risposta pi\`u precisa.
Il seguente teorema, sempre dovuto a Stanley, risolve il problema per la 
permutazione con la lunghezza pi\`u lunga.

\begin{theorem}[{\cite[Corollary~4.3]{stanley-84}}]
    Sia $\w_0$ la permutazione in $\sym_n$ che inverte l'ordine di $[n]$, ovvero
    $\w_0 = n, n-1, \dots , 1 $ in one-line notation.
    Allora vale
    \[
        \lvert \rw(\w_0) \rvert = f^{\lambda_0},
    \]
    dove $\lambda_0 = (n-1, n-2, \dots, 1)$ \`e la partizione il cui diagramma di
    Young \`e una ``scala'' con $\binom{n}{2}$ blocchi.
\end{theorem}

\begin{proof}
    Si verifica che il tableau $T_0$
    rappresentato in Figura~\ref{fig:tableau_scala} \`e un tableau ridotto per $\w_0$.
    Basta allora verificare che non ce ne siano altri, ovvero che se
    $T$ \`e un tableaux ridotto per $\w_0$, allora $T = T_0$.
    \begin{figure}[htbp]
        \centering
        \[
            % TODO fixa l'allineamento verticale qui
            T_0 = 
            \ytableaushort{
                {\scriptstyle n-1},
                {\scriptstyle n-2}{\scriptstyle n-1},
                {}{}{},
                3{}{}{},
                23{}{}{\scriptstyle n-1},
                123{}{\scriptstyle n-2}{\scriptstyle n-1}
                % 5,
                % 45,
                % 345,
                % 2345,
                % 12345
            }
        \]
        \caption{
            % Il tableau ridotto $T_0$ per $\w_0$, nel caso in cui $n = 6$.
            % La costruzione si generalizza per ogni $n$.
            Il tableau ridotto $T_0$ per la permutazione $\w_0 \in \sym_n$.
        }
        \label{fig:tableau_scala}
    \end{figure}

    Per ogni casella di $T$, consideriamo un percorso che parte dalla casella pi\`u
    in basso a sinistra e arriva a quella casella, muovendosi solo di una casella alla
    volta verso destra o verso l'alto.
    Ricordando l'Osservazione~\ref{rmk:tab_ridotto_crescenza},
    i numeri lungo questo percorso sono strettamente crescenti.

    Siccome $T$ contiene solo numeri da $1$ a $n-1$, il percorso sar\`a lungo al pi\`u
    $n-1$ caselle, da cui ogni casella di $T$ ha una distanza di Manhattan dalla 
    casella pi\`u in basso a sinistra minore o uguale a $n-1$.

    Le possibili caselle a distanza minore o uguale a $n-1$ sono esattamente 
    $\binom n 2$, ovvero quelle che compongono il diagramma ``a scala''.

    Poich\'e $\len(\w_0) = \binom n 2$, $T$ deve contenere $\binom n 2$ caselle,
    ovvero \emph{tutte} quelle a distanza minore o uguale a $n-1$, da cui 
    $\shape(T) = \lambda_0$.

    Per concludere che $T$ \`e uguale a $T_0$, basta osservare che $T_0$ \`e l'unico tableau 
    di quella forma con righe e colonne strettamente crescenti,
    che usa solo numeri da $1$ a $n-1$.
\end{proof}


    % Poich\'e $\len(\w_0) = \binom n 2$, $T$ deve contenere $\binom n 2$ caselle
    % a distanza minore o uguale a $n-1$, ovvero \emph{tutte}
    % quindi la sua forma \`e proprio quella ``a scala''.

    % Ricordiamo che Osservazione~\ref{rmk:tab_ridotto_crescenza} affermava che
    % righe e colonne di un tableau ridotto sono strettamente crescenti.
    % Usando questo fatto, ogni percorso come quello rappresentato sotto che si muove
    % solo verso destra e verso l'alto contiene numeri strettamente crescenti.
    % \[
    %     \ytableausetup{smalltableaux}
    %     \ytableaushort{
    %         \none,
    %         \none,
    %         \none      \none      {*(green)},
    %         \none      {*(green)} {*(green)},
    %         \none      {*(green)},
    %         {*(green)} {*(green)},
    %         {*(green)}
    %     }
    %     *{1,2,3,4,5,6,7}
    %     \ytableausetup{nosmalltableaux}
    % \]
    


\section{Estensione a fattorizzazioni crescenti}
    
In questa sezione estenderemo la corrispondenza EG a fattorizzazioni crescenti.

Ricordiamo che, dato $\w \in \sym_n$, $v = v^1 v^2 \cdots v^s$ \`e una 
\emph{fattorizzazione crescente} di $\w$ se $v$ \`e una parola ridotta per $\w$ e 
i fattori $v^i$ sono parole crescenti.

I singoli fattori $v^i$ possono anche essere la parola vuota.

Chiamiamo $\incrfact(\w)$ l'insieme (infinito) delle fattorizzazioni crescenti di $\w$.

Per ogni fattorizzazione crescente $v = v^1 v^2 \cdots v^s$ definisco il suo \emph{peso}
$\wt(v)$ come il vettore $\lambda$ che ha nella componente $i$-esima $\lambda_i$ la lunghezza
del fattore $v^i$.

Chiamo $\incrfact_\lambda(\w)$ l'insieme (questa volta finito) delle fattorizzazioni
crescenti di $\w$ con peso $\lambda$.

\begin{example}
    Per esempio, $v = 23|2||12|3$ \`e una fattorizzazione crescente di 
    $\w = s_2 s_3 s_2 s_1 s_2 s_3 = 4321 \in \sym_4$, con peso
    $\lambda = (2, 1, 0, 2, 1)$.
\end{example}

Il $P$-tableau di $v$ \`e ottenuto nello stesso modo, ovvero inserendo le lettere
di $v$ ignorando la suddivisione in fattori.

Invece, per ottenere il $Q$-tableau, scrivo $k$ in tutte le caselle ottenute dall'inserimento
delle lettere del fattore $v^k$ nel $P$-tableau.

Pi\`u formalmente, posso definire (in analogia a quanto definito prima) i $P$-tableaux 
intermedi $P^{(0)}$, $P^{(1)}$, \dots, $P^{(s)}$ come
\begin{align*}
    P^{(0)} &= \varnothing, \\
    P^{(k)} &= P^{(k-1)} \leftarrow v^k \qquad \text{se $k>0$};
\end{align*}
dove con $T \leftarrow v^k$ intendo il tableau ottenuto inserendo tutte le lettere di $v^k$
dentro $T$ una alla volta da sinistra a destra.

Come prima $P(v) = P^{(s)}$.
Come prima, se $( i, j)$ \`e una casella di $P^{(k)} \setminus P{(k-1)}$, allora scrivo
$k$ nella casella $(i,j)$ di $Q(v)$.

\begin{example}
    Se, come prima, $v = 23|2||12|3$, allora $P(v)$ e $Q(v)$ si ottengono come riportato:
    \begin{gather*}
        \ytableaunumber{2}
        \xrightarrow{\leftarrow 3}
        \ytableaunumber{2 3}
        \xrightarrow{\leftarrow 2}
        \ytableaunumber{3, 2 3}
        \xrightarrow{\leftarrow 1}
        \ytableaunumber{3, 2, 1 3}
        \xrightarrow{\leftarrow 2}
        \ytableaunumber{3, 2 3, 1 2}
        \xrightarrow{\leftarrow 3}
        \ytableaunumber{3, 2 3, 1 2 3}
        = P(v);
        \\
        \ytableaunumber{1}
        \longrightarrow
        \ytableaunumber{1 1}
        \longrightarrow
        \ytableaunumber{2, 1 1}
        \longrightarrow
        \ytableaunumber{4, 2, 1 1}
        \longrightarrow
        \ytableaunumber{4, 2 4, 1 1}
        \longrightarrow
        \ytableaunumber{4, 2 4, 1 1 5}
        =Q(v).
    \end{gather*}
    % TODO sistemare questo disegno
\end{example}

\begin{proposition}
    Per ogni una fattorizzazione crescente $v$, il tableau $Q(v)$ definito sopra \`e un 
    tableau semistandard.
\end{proposition}

\begin{proof}
    Basta verificare che $P^{(k)} \setminus P^{(k-1)}$ sia una \emph{striscia orizzontale}, 
    ovvero che non abbia due caselle sulla stessa colonna.
    Ma questo segue direttamente dalla crescenza del fattore $v^k$ e dal
    Lemma~\ref{lem:percorsi_strict} 
    sull'inserimento consecutivo di lettere crescenti.
\end{proof}

\begin{remark}
    Questa corrispondenza pi\`u generale per fattorizzazioni crescenti si riduce a quella
    vista in precedenza per parole ridotte scegliendo una fattorizzazione in cui ogni
    fattore ha una sola lettera.

    Per esempio, scegliendo $v = 2|3|2|1|2|3$, ritroviamo il $Q$-tableau calcolato
    nell'Esempio~\ref{ex:q_tab1}.
\end{remark}

\begin{theorem}
    \label{thm:eg2}
    % TODO definisci la roba in sto teorema
    Fissata una permutazione $\w \in \sym_n$ e un peso $\lambda \vDash \len(\w)$,
    la corrispondenza EG definita da
    \[
        \varphi_{\mathrm{EG}} \colon v \mapsto (P(v), Q(v))
    \]
    \`e una bigezione tra:
    \begin{itemize}
        \item l'insieme $\incrfact_\lambda(\w)$ delle fattorizzazioni crescenti
            di $\w$ con peso $\lambda$;
        \item l'insieme delle coppie $(P,Q)$ di tableaux di Young con la stessa forma,
            tali che $P$ sta in $\redtab(\w)$ e
            $Q$ \`e semistandard di tipo $\lambda$.
    \end{itemize}
\end{theorem}

\begin{proof}
    \`E vero per costruzione che $\wt(v) = \type(Q(v))$.

    Basta quindi verificare che l'algoritmo si pu\`o invertire, nonostante 
    $Q$ sia semistandard anzich\'e standard.
    Ma so gi\`a che le lettere di ogni fattore sono state inserite in ordine 
    crescente, quindi per il Lemma~\ref{lem:percorsi_strict} le caselle corrispondenti
    a quel fattore sono state aggiunte da quella pi\`u sinistra a quella pi\`u a destra.
    
    Allora, per invertire l'inserimento, si guarda la lettera pi\`u alta tra quelle che compaiono
    in $Q$ e tra le caselle in cui compare si sceglie quella pi\`u a destra, 
    Che \`e quella aggiunta per ultima.
    Poi si procede come nella dimostrazione del Teorema~\ref{thm:eg1}
    partendo da quella casella.
    Alla fine si suddivide la parola ottenuta in fattori, rispettando il tipo di $Q$.
\end{proof}

\begin{remark}
    % TODO o si toglie o si scrive
    Sostanzialmente sto standardizzando $Q$.
\end{remark}

Questo teorema \`e sufficiente per dare una scrittura delle funzioni simmetriche
di Stanley nella base delle funzioni simmetriche di Schur.

Ricordiamo che:
\begin{gather*}
    \stan_{\w^{-1}} = \sum_{ v \in \incrfact(\w) } x^{\wt(v)}, \\
    s_\lambda = \sum_{ T \in \ssyt(\lambda) } x^{\type(T)}.
\end{gather*}

\begin{corollary}
    Sia $\w \in \sym_n$ e sia $\stan_{\w^{-1}}$ la funzione simmetrica di Stanley associata
    a $\w^{-1}$.
    Allora vale:
    % \[
    %     \stan_{\w^{-1}} = \sum_{T \in \redtab(w)} s_{\shape(T)},
    % \]
    % o, equivalentemente:
    \[
        \stan_{\w^{-1}} = \sum_{\lambda} a_\w^\lambda \, s_\lambda,
    \]
    dove $a_\w^\lambda$ \`e il numero di tableaux semistandard il cui row-reading
    \`e una parola ridotta per $\w$.
    In simboli:
    \begin{equation}
        \label{eq:def_a_w_lambda}
        a_\w^\lambda \coloneq \lvert \left\{ 
            T \in \redtab(\w) \mid \shape(T) = \lambda
        \right\} \rvert .
    \end{equation}
    
\end{corollary}

\begin{proof}
    % TODO fai proof (solo conti)
\end{proof}


\section{Struttura di cristallo}

In questa sezione vediamo come la corrispondenza di Edelman-Greene sia anche
un isomorfismo di cristalli, ovvero la struttura di cristallo sulle fattorizzazioni
crescenti definita nel capitolo~\ref{ch:stanley_symm} viene mappata nella struttura 
di cristallo sui tableaux semistandard definita del capitolo~\ref{ch:crystals}.

Fissata una permutazione $\w \in \sym_n$, e un intero positivo $\ell$ che 
assumeremo essere ``sufficientemente grande'',
la corrispondenza di Edelman-Greene pu\`o
essere vista come un bigezione tra l'insieme $\incrfact^\ell(\w)$ delle
fattorizzazioni crescenti di $\w$ con $\ell$ fattori e
l'unione disgiunta di opportuni insiemi di tableaux semistandard,
indicizzata dai tableaux ridotti  di $\w$.

Pi\`u precisamente:
\[
    \varphi_\mathrm{EG} \colon 
    \incrfact^\ell(\w) \to
    \bigsqcup_{P \in \redtab(\w)}
    % \!\!\! NOTE pondere se mettere meno spazio
    \ssyt^\ell\bigl(\shape(P)\bigr),
\]
dove per ogni partizione $\lambda$,
$\ssyt^\ell(\lambda)$ \`e l'insieme dei tableaux semistandard di forma $\lambda$ in 
cui compaiono solo lettere da $1$ a $\ell$.

Ricordiamo che nel Capitolo~\ref{ch:stanley_symm}
abbiamo definito una struttura di cristallo $B(\w)$ sull'insieme 
$\incrfact^\ell(\w)$.

Ricordiamo anche che nella Sezione~\ref{sec:crys_tableaux} abbiamo definito
una struttura di cristallo $B(\lambda)$ sull'insieme $\ssyt^\ell(\lambda)$,
quindi possiamo vedere il codominio come la somma diretta di questi cristalli:
\[
    \bigoplus_{P \in \redtab(\w)} B\bigl(\shape(P)\bigr)
    =
    \bigoplus_\lambda B(\lambda)^{\oplus a_\w^\lambda},
\]
dove $a_\w^\lambda$, definito in \eqref{eq:def_a_w_lambda}, \`e il numero di 
tableaux semistandard il cui row-reading \`e una parola ridotta per $\w$.

Il fatto sorprendente \`e che la corrispondenza EG tra questi due insiemi, oltre a
essere una bigezione che preserva il peso, \`e anche un isomorfismo di cristalli.

\begin{theorem}
    [{\cite[Theorem~4.11]{morse-schilling-15}}]
    \label{thm:eg_isomorfismo}
    Per ogni permutazione $\w \in \sym_n$, la mappa di Edelman-Greene
    \[
        \varphi_\mathrm{EG} \colon 
        B(\w) \to 
        \bigoplus_\lambda B(\lambda)^{\oplus a_\w^\lambda}
    \]
    \`e un isomorfismo di cristalli.

    In particolare:
    \begin{enumerate}
        \item l'azione di $f_i$ ed $e_i$ sulla fattorizzazione crescente $v$, se definita,
            non cambia il suo $P$-tableau, ovvero
            \begin{align*}
                P(f_i(v)) &= P(v) && \text{se $f_i(v) \neq \zero$,} \\
                P(e_i(v)) &= P(v) && \text{se $e_i(v) \neq \zero$;}
            \end{align*}
        \item l'azione di $f_i$ ed $e_i$ sulla fattorizzazione crescente $v$ commuta
            con quella sul suo $Q$-tableau, ovvero
            \begin{align*}
                Q(f_i(v)) = f_i(Q(v)), \\
                Q(e_i(v)) = e_i(Q(v)).
            \end{align*}
    \end{enumerate}





\end{theorem}

Prima di dimostrare questo teorema, introduciamo la nozione di
\emph{Coxeter-Knuth equivalenza}, un raffinamento della Coxeter equivalenza introdotto
da Edelman e Greene in \cite[Definition~6.19]{edelman-greene-87}.

\begin{definition}
    Due parole ridotte $v$ e $v'$ si dicono \emph{Coxeter-Knuth equivalenti}
    (e lo scrivo $v \approx v'$) se si possono ottenere
    una dall'altra tramite una serie di operazioni di Coxeter-Knuth, ovvero
    sostituzioni della forma
    \begin{subequations}
        \begin{align}
            \label{subeq:ck_braid}
            x(x+1)x &\leftrightarrow (x+1)x(x+1),  \\
            \label{subeq:ck_sx}
            y x z  &\leftrightarrow y z x  && \text{se $x<y<z$ e $\lvert x-z \rvert \geq 2$, } \\
            \label{subeq:ck_dx}
            x z y  &\leftrightarrow z x y  && \text{se $x<y<z$ e $\lvert x-z \rvert \geq 2$. }
        \end{align}
    \end{subequations}
\end{definition}

In sostanza, posso scambiare due lettere che commutano $x$ e $z$ solo se,
da uno dei due lati, c'\`e una lettera $y$ numericamente compresa tra le due.

\begin{lemma}[{\cite[Lemma~6.23]{edelman-greene-87}}]
    \label{lem:ptab_ck_equiv}
    Ogni parola ridotta $v$ \`e Coxeter-Knuth equivalente al row-reading del suo
    $P$-tableau, in simboli:
    \[
        v \approx \rowreading(P(v)).
    \]
\end{lemma}

\begin{proof}
    Ripercorrendo i passaggi dell'Osservazione~\ref{rmk:coxeter_righe}, si osserva
    che tutte le operazioni di Coxeter utilizzate sono gi\`a operazioni di Coxeter-Knuth.
\end{proof}

\begin{lemma}
    [{\cite[Theorem~6.24]{edelman-greene-87}}]
    \label{lem:ptab_ck_equiv_doppia_freccia}
    Due parole ridotte $v$ e $v'$ sono Coxeter-Knuth equivalenti se e solo se
    i loro $P$-tableaux sono uguali. In simboli:
    \[
        v \approx v' \iff P(v) = P(v').
    \]
\end{lemma}

\begin{proof}
    Se $P(v) = P(v')$, allora, per il Lemma~\ref{lem:ptab_ck_equiv},
    \[
        v \approx P(v) = P(v') \approx v'.
    \]

    Per l'altra implicazione, che si riconduce a una lunga casistica,
    rimandiamo alla dimostrazione di Edelman e Greene in
    \cite[Theorem~6.4]{edelman-greene-87}.
\end{proof}

Ritorniamo ora al cristallo di Morse e Schilling $B(\w)$ per ogni $\w \in \sym_n$.

\begin{lemma}
    \label{lem:ms_f_ck_equivalenza}
    Per ogni $i \in I$ e per ogni fattorizzazione crescente $v \in B(\w)$
    tale che $f_i(v) \neq \zero$,
    le parole ridotte associate a $v$ e a $f_i(v)$ sono Coxeter-Knuth equivalenti.
\end{lemma}

\begin{proof}
    Ripercorriamo la dimostrazione della
    Proposizione~\ref{prop:ms_f_coxeter_equivalenza}, che dimostrava la
    Coxeter equivalenza tra $v$ e $f_i(v)$.
    Si osserva, anche grazie al Lemma~\ref{lem:divisione_in_3}, che tutte le operazioni
    di Coxeter utilizzate sono gi\`a operazioni di Coxeter-Knuth.

    Per le propriet\`a dell'operatore $f_i$ descritte nella
    Sezione~\ref{sec:cristallo_morse_schilling}, basta guardare il caso con due fattori.
    Utilizziamo inoltre il pairing della Definizione~\ref{def:ms_i_pairing} e la suddivisione
    del Lemma~\ref{lem:divisione_in_3}.
    
    Mostriamo come ottenere $f_1(v)$ da $v = v^1v^2$, utilizzando come esempio la parola
    ridotta dell'Esempio~\ref{ex:f_agisce_su_incrfact}, dove le lettere in 
    grassetto sono quelle non appaiate:
    \[
        v = v^1  v^2 = 
        ( \mathbf{1}, 2, \mathbf{3}, 9, 11 \mid \mathbf{12}, 13, 14 \mid 18 ) \,
        ( 1, 5, 8 \mid 12, 13 \mid \mathbf{15}, 16).
    \]

    % TODO: emoji del teschio + rifai con lem:divisione_sgravo_comparativo
    
    Per prima cosa sposto a sinistra le lettere del primo blocco del
    secondo fattore.
    Per il Lemma~\ref{lem:lett_appaiate_ordinate}, queste lettere
    sono ordinatamente minori delle lettere appaiate nel primo blocco del primo
    fattore.
    Allora posso muoverle (una alla volta) verso sinistra fino a raggiungere
    ciascuna la propria lettera corrispondente (che pu\`o anche non essere quella
    con cui erano appaiate), ottenendo:
    \[
        ( \mathbf{1}, 2 \text{-} 1, \mathbf{3}, 9 \text{-} 5,
        11 \text{-} 8\mid \mathbf{12}, 13, 14 \mid 18 ) \,
        ( 12, 13 \mid \mathbf{15}, 16).
    \]
    Per ogni lettera, la prima operazione di Coxeter-Knuth \`e della forma in
    \eqref{subeq:ck_dx}, tutte le successive sono della forma in 
    \eqref{subeq:ck_sx}.

    Nello stesso modo, tutte le lettere del terzo blocco del primo fattore si
    possono spostare a destra, fino alla lettera ordinatamente corrispondente tra
    quelle appaiate, ottenendo:
    \[
        ( \mathbf{1}, 2 \text{-} 1, \mathbf{3}, 9 \text{-} 5,
        11 \text{-} 8\mid \mathbf{12}, 13, 14  ) \,
        ( 12, 13 \mid \mathbf{15}, 18 \text{-} 16).
    \]
    Si verifica che tutte le operazioni di Coxeter richieste dal Lemma~\ref{lem:parole_equiv}
    sono gi\`a operazioni di Coxeter-Knuth, quindi ottengo:
    \[
        ( \mathbf{1}, 2 \text{-} 1, \mathbf{3}, 9 \text{-} 5,
        11 \text{-} 8\mid  13, 14  ) \,
        ( 12, 13, \mathbf{14} \mid \mathbf{15}, 18 \text{-} 16).
    \]
    Alla fine, tutte le lettere spostate durante i primi step
    possono ritornare alle loro posizioni originali, ottenendo
    \[
        f_1(v) = 
        ( \mathbf{1}, 2, \mathbf{3}, 9, 11 \mid 13, 14 \mid 18 ) \,
        ( 1, 5, 8 \mid 12, 13, \mathbf{14} \mid \mathbf{15}, 16).
        \qedhere
    \]
\end{proof}

Possiamo ora procedere con la dimostrazione del Teorema~\ref{thm:eg_isomorfismo}.

\begin{proof}
    % [Dimostrazione del Teorema~\ref{thm:eg_isomorfismo}]
    Prima dimostriamo il punto~??1.
    Per il Lemma~\ref{lem:ms_f_ck_equivalenza}, le parole ridotte $v$ e $f_i(v)$ sono 
    Coxeter-Knuth equivalenti e, per il Lemma~\ref{lem:ptab_ck_equiv_doppia_freccia},
    parole Coxeter-Knuth equivalenti hanno lo stesso $P$-tableau e questo conclude
    il punto~??1.

    Prima di dimostrare il punto~??2, introduciamo un po' di notazione.

    Denotiamo le lettere dello $i$-esimo e dello $i+1$-esimo fattore di $v$ con
    \begin{align*}
        v^i &= x_1 x_2 \cdots x_{\alpha_i}, \\
        v^{i+1} &= y_1 y_2 \cdots y_{\alpha_{i+1}}.
    \end{align*}
    Sia inoltre $x_j$ l'ultima lettera non appaiata di $v^i$, ovvero quella 
    che verr\`a rimossa dall'azione di $f_i$ e sostituita con un'opportuna lettera
    in $v^{i+1}$.
    Sia inoltre $y_k$ la pi\`u grande lettera di $v^{i+1}$ minore di $x_j$,
    ovvero l'ultima lettera del primo blocco secondo la suddivisione del 
    Lemma~\ref{lem:divisione_in_3}.


    In questa dimostrazione useremo pi\`u volte il Lemma~\ref{lem:percorsi_weak}
    con maggiore generalit\`a di quella dell'enunciato originale.
    Infatti il lemma rimane valido anche se la seconda lettera non \`e la lettera
    inserita immediatamente dopo, a patto che il percorso d'inserimento della prima
    lettera non sia stato ``toccato'' dall'inserimento di altre lettere intermedie.

    Mostriamo subito questo ragionamento in azione.

    \begin{lemma}
        \label{lem:thm_trascuro_inizio}
        I percorsi di inserimento delle lettere $y_1, \dots, y_k$ sono tutti strettamente
        a sinistra del percorso di inserimento di $x_j$.
    \end{lemma}

    \begin{proof}[Dimostrazione del lemma]
        Per il Lemma~\ref{lem:divisione_sgravo_comparativo}, la lettera $y_1$
        \`e minore della lettera $x_{j-k}$.
        I percorsi d'inserimento delle lettere di $v^i$ a destra di $x_{j-k}$ sono strettamente
        a destra del percorso di $x_{j-k}$ per il Lemma~\ref{lem:percorsi_strict},
        quindi non lo hanno ``toccato'' e il percorso di inserimento di $y_1$ sar\`a
        debolmente a sinistra di quello di $x_{j-k}$.

        Successivamente, devo inserire la lettera $y_2 $, che per lo stesso lemma
        \`e minore della lettera $x_{j-k+1}$.
        Il percorso d'inserimento di $x_{j-k+1}$ non pu\`o essere stato ``toccato'' n\'e dalle
        successive lettere di $v^i$, n\'e da $y_1$, in quanto il suo percorso \`e a sinistra
        di quello di $x_{j-k}$.
        Allora concludo che il percorso d'inserimento di $y_2$ \`e debolmente a sinistra
        di quello di $x_{j-k+1}$.

        Lo stesso ragionamento va rifatto per ogni lettera fino a $y_k$, il cui percorso
        d'inserimento sar\`a debolmente a sinistra di quello di $x_{j-1}$.
        Ma allora tutti i percorsi sono, in particolare, strettamente a sinistra
        del percorso di $x_j$.
        % TODO ponderare se mettere esempio qui.
    \end{proof}

    Studiamo ora i percorsi di inserimento delle lettere di $v^i$ e come cambiano
    rimuovendo $x_j$.

    Per il Lemma~\ref{lem:percorsi_strict}, per ogni $r = 1, \dots, \alpha_i$
    il percorso di inserimento di $x_{r+1}$ \`e disgiunto e strettamente
    a destra del percorso di $x_r$.

    Rimuovendo $x_j$, chiaramente i percorsi di $x_1, \dots, x_{j-1}$ non cambiano.
    Inserendo $x_{j+1}$ senza aver inserito $x_j$, ci sono due situazioni possibili:
    \begin{itemize}
        \item la lettera $a$ che avrebbe spinto $x_j$ \`e anche maggiore di $x_{j+1}$,
            allora $x_{j+1}$ spinge $a$ e devia proseguendo lungo tutto il percorso
            di inserimento originario di $x_j$;
        \item la lettera $a$ \`e troppo piccola per essere spinta da $x_{j+1}$,
            allora $x_{j+1}$ spinge la stessa lettera che avrebbe spinto se ci fosse stato
            $x_j$ e fa un passo lungo il proprio percorso di inserimento originale.
    \end{itemize}
    Inoltre, finch\'e non si verifica la prima situazione,
    lo stesso confronto tra le lettere spinte nei due percorsi di inserimento originali
    si ripete ad ogni riga.
    Alla fine, il nuovo percorso di inserimento di $x_{j+1}$ \`e composto alla
    base da un pezzo del proprio percorso originale fino a un punto in cui (eventualmente)
    devia sul percorso originale di $x_j$.

    Inserendo $x_{j+2}$ (e analogamente le lettere successive), il percorso prosegue
    sul vecchio tragitto finch\'e $x_{j+1}$ non \`e ancora stato deviato,
    poi incontra lo stesso confronto descritto sopra che determiner\`a se
    il percorso prosegue sul vecchio tragitto o se devia sul tragitto originale
    di $x_{j+1}$, da lui non pi\`u utilizzato.

    \begin{example}
        \label{ex:thm_tab_con_inspath}
        In Figura~\ref{fig:tab_con_insertionp} mostriamo l'inserimento nel tableau $T$ delle lettere
        \[
            x_1 x_2 x_3 x_4 = 4, 6, 10, 16
        \]
        e successivamente l'inserimento escludendo 
        la lettara $x_1$.
    \end{example}


    % TODO pondera dove mettere la figura.
    \begin{figure}[htbp]
        \centering
        \[
            \ytableausetup{aligntableaux=center}
            T = \ytableaunumber{
                { 8} {15} {19},
                { 6} {13} {14} {18} {22},
                { 2} { 3} { 9} {15} {16} {17} {24},
                { 1} { 2} { 5} { 7} { 8} {12} {13} {18}
            }
            \ytableausetup{aligntableaux=bottom}
        \]

        \begin{tikzpicture}[remember picture]
            \node (T1) {
                    \ytableaunumber{
                        {\tmn{A51}{15}} {\tmn{A52}{19}},
                        {\tmn{A41}{ 8}} {\tmn{A42}{13}} {\tmn{A43}{18}} {\tmn{A44}{22}},
                        {\tmn{A31}{ 6}} {\tmn{A32}{ 9}} {\tmn{A33}{14}} {\tmn{A34}{15}} {\tmn{A35}{16}} {\tmn{A36}{24}},
                        {\tmn{A21}{ 2}} {\tmn{A22}{ 3}} {\tmn{A23}{ 5}} {\tmn{A24}{ 7}} {\tmn{A25}{12}} {\tmn{A26}{17}} {\tmn{A27}{18}},
                        {\tmn{A11}{ 1}} {\tmn{A12}{ 2}} {\tmn{A13}{ 4}} {\tmn{A14}{ 6}} {\tmn{A15}{ 8}} {\tmn{A16}{10}} {\tmn{A17}{13}} {\tmn{A18}{16}}
                    }
                };
            \node (T2) at (6,0) {
                    \ytableaunumber{
                        {\tmn{B51}{15}} {\tmn{B52}{19}},
                        {\tmn{B41}{ 8}} {\tmn{B42}{13}} {\tmn{B43}{18}} {\none[\bullet]},
                        {\tmn{B31}{ 6}} {\tmn{B32}{ 9}} {\tmn{B33}{14}} {\tmn{B34}{15}} {\tmn{B35}{22}} {\tmn{B36}{14}},
                        {\tmn{B21}{ 2}} {\tmn{B22}{ 3}} {\tmn{B23}{ 7}} {\tmn{B24}{12}} {\tmn{B25}{16}} {\tmn{B26}{17}} {\tmn{B27}{18}},
                        {\tmn{B11}{ 1}} {\tmn{B12}{ 2}} {\tmn{B13}{ 5}} {\tmn{B14}{ 6}} {\tmn{B15}{ 8}} {\tmn{B16}{10}} {\tmn{B17}{13}} {\tmn{B18}{16}}
                    }
                };

            \begin{scope}[on background layer, overlay]
                \draw[green, opacity=0.7, line width=11pt, line cap=round, line join=round] 
                    (A13.center) to
                    (A23.center) to
                    (A32.center) to
                    (A42.center) to
                    (A51.center);
                \draw[green, opacity=0.7, line width=11pt, line cap=round, line join=round] 
                    (A14.center) to
                    (A24.center) to
                    (A34.center) to
                    (A43.center) to
                    (A52.center);
                \draw[green, opacity=0.7, line width=11pt, line cap=round, line join=round] 
                    (A16.center) to
                    (A25.center) to
                    (A35.center) to
                    (A44.center);
                \draw[green, opacity=0.7, line width=11pt, line cap=round, line join=round] 
                    (A18.center) to
                    (A27.center) to
                    (A36.center);

                \draw[green, opacity=0.7, line width=11pt, line cap=round, line join=round] 
                    (B14.center) to
                    (B23.center) to
                    (B32.center) to
                    (B42.center) to
                    (B51.center);
                \draw[green, opacity=0.7, line width=11pt, line cap=round, line join=round] 
                    (B16.center) to
                    (B24.center) to
                    (B34.center) to
                    (B43.center) to
                    (B52.center);
                \draw[green, opacity=0.7, line width=11pt, line cap=round, line join=round] 
                    (B18.center) to
                    (B27.center) to
                    (B36.center);

                % \node[circle, opacity=0.7, draw=blue, ultra thick, inner sep=5pt]
                %     (bucaA1) at (A32.center) {};
                % \node[circle, opacity=0.7, draw=blue, ultra thick, inner sep=5pt]
                %     (bucaB1) at (A24.center) {};
                % \draw[blue, opacity=0.7, ultra thick]
                %     (bucaA1) to (bucaB1);
                % \node[circle, opacity=0.7, draw=blue, ultra thick, inner sep=5pt]
                %     (bucaA1) at (A34.center) {};
                % \node[circle, opacity=0.7, draw=blue, ultra thick, inner sep=5pt]
                %     (bucaB1) at (A25.center) {};
                % \draw[blue, opacity=0.7, ultra thick]
                %     (bucaA1) to (bucaB1);

            \end{scope}
        \end{tikzpicture}
        \caption{L'inserimento nel tableau $T$ prima delle lettere $4$, $6$, $10$, $16$,
            poi solo delle lettere $6$, $10$, $16$.
            I percorsi di inserimento sono evidenziati.
        }
        \label{fig:tab_con_insertionp}
    \end{figure}

    Con quanto detto abbiamo dimostrato che, rimuovendo $x_j$, i percorsi di inserimento
    di $x_{j+1}, \dots, x_{\alpha_i}$ o deviano a un certo punto sul percorso alla propria
    sinistra, oppure proseguono sul proprio vecchio percorso.

    Questo significa che la forma del $P$-tableau $P^{(i)}(f_i(v))$
    (quello senza $x_j$) si ottiene da quella di $P^{(i)}(v)$ (quello con $x_j$)
    rimuovendo una casella, in particolare la punta originale dell'ultimo percorso
    deviato.
    La casella rimossa \`e indicata con un pallino nella Figura~\ref{fig:tab_con_insertionp}.

    Poich\'e, per la prima parte del teorema, il $P$-tableau finale deve essere uguale,
    la casella rimossa deve venire aggiunta dal fattore $v^{i+1}$ e il $Q$-tableau
    di $f_i(v)$ si ottiene da quello di $v$ cambiando da $i$ a $i+1$ la lettera 
    nella casella mancante, ovvero la punta originale dell'ultimo percorso che devia.
    Il resto della dimostrazione serve per verificare che la lettera $i$ che viene
    cambiata in $i+1$ sia proprio l'\emph{ultima} lettera $i$ \emph{spaiata} nel 
    row-reading del $Q$-tableau in cui le lettere $i$ sono sostituite da perentesi
    chiuse ``$)$'' e le lettere $i+1$ da parentesi aperte ``$\lparen$'',
    come richiesto dalla definizione di cristallo di tableaux data nella 
    Sezione~\ref{sec:crys_tableaux}.

    Studiamo allora pi\`u precisamente dove avviene la deviazione sul percorso alla
    sinistra.

    Chiamiamo $a_1, a_2, \dots$ le lettere che leggiamo lungo il percorso di inserimento
    di $x_j$ e similmente chiamiamo $b_1, b_2, \dots$ le lettere lungo il percorso di $x_{j+1}$,
    dove per convezione $a_r = \infty$ se il percorso non raggiunge la riga $r$.

    Nel seguito dimostriamo che la deviazione avviene 
    alla prima riga $r$ in cui $a_{r+1} > b_r$.

    Assumiamo prima che non avvengano inserimenti speciali.
    Prima degli inserimenti, la lettera $a_{r+1}$ si trovava nella riga $r$
    assieme a $b_{r+1}$.
    Quando inserisco $x_{j+1}$ senza aver prima inserito $x_j$, la lettera $b_r$, che
    avrebbe altrimenti spinto $b_{r+1}$, spinger\`a invece $a_{r+1}$ in quanto
    % Quando non viene inserito $x_j$, la lettera $a_{r+1}$ si trova nella riga $r$.
    % Allora, la lettera $b_r$ che aveva spinto $b_{r+1}$ spinger\`a invece
    % $a_{r+1}$ in quanto
    \[
        b_r < a_{r+1} < b_{r+1}.
    \]

    Per via degli inserimenti speciali, pu\`o succedere che la lettera inserita
    in una riga sia la lettera che leggiamo nel percorso \emph{diminuita di uno}.
    Bisogna quindi ripetere il ragionamento sostituendo alcune delle lettere
    con $b_r-1$, $a_{r+1}-1$ o $b_{r+1}-1$.
    i vari casi portano allo stesso risultato oppure ad assurdi.

    Anche per mostrare il viceversa, (ovvero se $a_{r+1} \leq b_r$ allora
    non c'\`e deviazione), le eventuali diminuzioni delle lettere non causano
    problemi nella dimostrazione (che \`e immediata nel caso senza inserimenti
    speciali).
    % NOTE boia, questo proprio non \`e stato dimostrato. Pondera se aggiungere
    % un esempio,  :,(

    Chiaramente lo stesso discorso si pu\`o applicare anche per 
    i percorsi di inserimento successivi a quello di $x_{j+1}$ e quindi caratterizza
    tutti i punti in cui un percorso devia sul precedente.
    Per il resto della dimostrazione, mi riferisco a questi punti di deviazione come
    \emph{buche}.

    Nella Figura~\ref{fig:buche_esempi} sono raffigurate le due buche 
    dell'Esempio~\ref{ex:thm_tab_con_inspath}, le lettere cerchiate sono quelle
    tra cui si deve verificare la disuguaglianza.

    \begin{figure}
        \centering
        \begin{tikzpicture}[remember picture]
            \node (esempi) {
                    \ytableaushort{
                        {\none[\tmn{A14}{}]},
                        {\tmn{A13}{a_2}} {\none[\cdots]} \none {\none[\tmn{B13}{}]},
                        \none \none {\tmn{A12}{a_1}} {\tmn{B12}{b}},
                        \none \none {\none[\tmn{A11}{}]} {\none[\tmn{B11}{}]}
                    }
                    \qquad
                    \qquad
                    \ytableaunumber{
                        {\none[\tmn{A24}{}]},
                        {\tmn{A23}{9}} \none {\none[\tmn{B23}{}]},
                        \none {\tmn{A22}{5}} {\tmn{B22}{7}},
                        \none {\none[\tmn{A21}{}]} {\none[\tmn{B21}{}]}
                    }
                    \qquad
                    \qquad
                    \ytableaunumber{
                        {\none[\tmn{A34}{}]},
                        {\tmn{A33}{15}} {\none[\tmn{B33}{}]},
                        {\tmn{A32}{7}} {\tmn{B32}{12}},
                        {\none[\tmn{A31}{}]} {\none[\tmn{B31}{}]}
                    }
                };

            \begin{scope}[on background layer, overlay]
                \draw[green, opacity=0.7, line width=11pt, line cap=round, line join=round] 
                    (A11.center) to
                    (A12.center) to
                    (A13.center) to
                    (A14.center);
                \draw[green, opacity=0.7, line width=11pt, line cap=round, line join=round] 
                    (B11.center) to
                    (B12.center) to
                    (B13.center);
                \draw[green, opacity=0.7, line width=11pt, line cap=round, line join=round] 
                    (A21.center) to
                    (A22.center) to
                    (A23.center) to
                    (A24.center);
                \draw[green, opacity=0.7, line width=11pt, line cap=round, line join=round] 
                    (B21.center) to
                    (B22.center) to
                    (B23.center);
                \draw[green, opacity=0.7, line width=11pt, line cap=round, line join=round] 
                    (A31.center) to
                    (A32.center) to
                    (A33.center) to
                    (A34.center);
                \draw[green, opacity=0.7, line width=11pt, line cap=round, line join=round] 
                    (B31.center) to
                    (B32.center) to
                    (B33.center);

                \node[circle, opacity=0.7, draw=blue, ultra thick, inner sep=5pt]
                    (bucaA1) at (A13.center) {};
                \node[circle, opacity=0.7, draw=blue, ultra thick, inner sep=5pt]
                    (bucaB1) at (B12.center) {};
                \draw[blue, opacity=0.7, ultra thick]
                    (bucaA1) to (bucaB1);
                \node[circle, opacity=0.7, draw=blue, ultra thick, inner sep=5pt]
                    (bucaA2) at (A23.center) {};
                \node[circle, opacity=0.7, draw=blue, ultra thick, inner sep=5pt]
                    (bucaB2) at (B22.center) {};
                \draw[blue, opacity=0.7, ultra thick]
                    (bucaA2) to (bucaB2);
                \node[circle, opacity=0.7, draw=blue, ultra thick, inner sep=5pt]
                    (bucaA3) at (A33.center) {};
                \node[circle, opacity=0.7, draw=blue, ultra thick, inner sep=5pt]
                    (bucaB3) at (B32.center) {};
                \draw[blue, opacity=0.7, ultra thick]
                    (bucaA3) to (bucaB3);
            \end{scope}
        \end{tikzpicture}

        \caption{La struttura generica di una buca, con $a_2 > b$, e le due
        buche dell'esempio nella Figura~\ref{fig:tab_con_insertionp}.}
        \label{fig:buche_esempi}
    \end{figure}

    % NOTE probabilmente questo non \`e inutile
    Nel seguito dimostriamo che tutte le buche hanno la struttura descritta
    nella Figura~\ref{fig:buche_esempi}, ovvero le caselle alla base, contenenti
    $a_1$ e $b$, sono contigue, cio\`e non ci sono altre caselle in mezzo, e
    i due cammini di inserimento si toccano in quel punto.

    Se le caselle di $a_1$ e di $b$ non fossero contigue,
    ci sarebbe una lettera $c$, non toccata dagli inserimenti, con $a_1 < c < b$.
    Inoltre, poich\'e la casella di $a_1$ era precedentemente occupata da $a_2$
    (o $a_2-1$), vale anche $a_2 \leq c < b$, che \`e assurdo perch\'e
    in una buca vale $b < a_2$.

    Ora dobbiamo studiare cosa succede inserendo le lettere 
    $y_{k+1}, y_{k+2}, \dots, y_{\alpha_{i+1}}$, sempre con l'obbiettivo di mostrare
    che la casella che viene rimossa quando non si inserisce $x_j$ \`e l'ultima
    parentesi chiusa non appaiata quando scrivo una ``$)$'' all'estremit\`a dei
    percorsi d'inserimento di $v^i$ e una ``$\lparen$'' all'estremit\`a di quelli
    di $v^{i+1}$.

    Il Lemma~\ref{lem:thm_trascuro_inizio} ci assicura che l'inserimento delle precedenti
    lettere di $v^{i+1}$ non ha mai toccato il percorso di inserimento di $x_j$ e
    le caselle alla sua destra.

    Anche per il Lemma~\ref{lem:divisione_sgravo_comparativo}, abbiamo che 
    $x_j \leq y_{k+1} < x_{j+1}$.
    Nel seguito voglio dimostrare che:
    \begin{itemize}
        \item se c'\`e una buca tra i percorsi di $x_j$ e di $x_{j+1}$, allora il percorso
            di $y_{k+1}$ scavalca quello di $x_j$ e finisce alla sua sinistra;
        \item se non c'\`e alcuna buca tra i due percorsi, allora il percorso di
            $y_{k+1}$ rimane strettamente a destra di quello di $x_j$.
    \end{itemize}

    Per il Lemma~\ref{lem:percorsi_weak}, il percorso di $y_{k+1}$ rimane debolmente
    a sinistra di quello di $x_{j+1}$.



    
    

\end{proof}

